\subsection{Lebesgue Integration}

% **Role:** modern integration and almost-everywhere analysis.
%
% Nonnegative integration, integrable functions, comparison with Riemann integration, absolute integrability.

Lebesgue integration is calibrated on indicators, extended linearly to nonnegative simple functions, then extended by monotone approximation and positive--negative decomposition.

\begin{perspective} \label{per:why-general-integral-uses-simple-minorants} \hlt{Simple Minorants Are the Right Primitive}\\
For Lebesgue measure on $\R$, bounded measurable functions supported on finite-measure sets form a useful intermediate class; arbitrary measure spaces need not admit that reduction.\\
Let $X\neq\emptyset$, $\mathcal{M}=\{\emptyset,X\}$, and $\mu(X)=+\infty$. If $g$ is measurable and vanishes outside some $F\in\mathcal{M}$ with $\mu(F)<+\infty$, then $F=\emptyset$ and $g=0$; hence finite-measure-support minorants of $\mathbf{1}_X$ have supremal integral $0$, while the required normalization is $\int_X\mathbf{1}_X\,d\mu=\mu(X)=+\infty$. Thus the defining minorants must include $\mathbf{1}_A$ $\forall A\in\mathcal{M}$, and more generally all nonnegative simple $s\leq f$.\\
\hlt{Key Insight}: The primitive calibration is $\int_X\mathbf{1}_A\,d\mu=\mu(A)$, with $\mu(A)=+\infty$ allowed.
\end{perspective}

\begin{definition} \label{def:simple-functions-and-lebesgue-integral} \hlt{Lebesgue Integral}\\
Let $(X,\mathcal{M},\mu)$ be a measure space and $E\in\mathcal{M}$; let $\mathcal{S}^+$ be nonnegative measurable simple functions on $X$.\\
Write $s\in\mathcal{S}^+$ as $s=\sum\limits_{i=1}^{n}c_i\mathbf{1}_{A_i}$ with $c_i\geq0$ and pairwise disjoint $A_i\in\mathcal{M}$.\\
For measurable $f\geq0$, put $\mathcal{S}^+(f)\defeq\{s\in\mathcal{S}^+:s\leq f\}$; for measurable $f:X\to\overline{\R}$, put $f^\pm\defeq(\pm f)\vee0$.
\begin{enumerate}[label=(\roman*)]
\setlength{\itemsep}{0pt}
\item \hlt{Simple}: $\mathcal{I}_E(s)\defeq\sum\limits_{i=1}^{n}c_i\mu(E\cap A_i)$, with $0\cdot(+\infty)\defeq0$.
\item \hlt{Nonnegative}: $\int_Ef\,d\mu\defeq\sup\limits_{s\in\mathcal{S}^+(f)}\mathcal{I}_E(s)$; $f$ integrable on $E\Longleftrightarrow\int_Ef\,d\mu<+\infty$.
\item \hlt{Signed}: if $\min\{\int_Ef^+\,d\mu,\int_Ef^-\,d\mu\}<+\infty$, set $\int_Ef\,d\mu\defeq\int_Ef^+\,d\mu-\int_Ef^-\,d\mu$.
\item \hlt{$L^1$}: $f\in\mathscr{L}^1(\mu)$ on $E$ iff $f$ is measurable and $\int_E\abs f\,d\mu<+\infty$.
\end{enumerate}
When $E=X$, write $\int f\,d\mu\defeq\int_Xf\,d\mu$ and drop ``on $E$''.
\end{definition}

\begin{proposition} \label{prop:signed-integral-independent-of-decomposition} \hlt{Signed Decompositions Give the Same Integral}\\
Let $(X,\mathcal{M},\mu)$ be a measure space, let $E\in\mathcal{M}$, and let $f,f_1,f_2:X\to\overline{\R}$ be measurable.\\
Assume $f_i\geq0$, $\int_Ef_i\,d\mu<+\infty$ for $i=1,2$, and $f=f_1-f_2$ on $E$. Then:
\begin{enumerate}[label=(\roman*)]
\setlength{\itemsep}{0pt}
\item $f\in\mathscr{L}^1(\mu)$ on $E$.
\item $\int_E f\,d\mu=\int_E f_1\,d\mu-\int_E f_2\,d\mu$.
\end{enumerate}
\begin{proof}
\begin{enumerate}[label=(\roman*)]
\setlength{\itemsep}{0pt}
\item Since $-f_2\leq f\leq f_1$ on $E$, $f^+\leq f_1$ and $f^-\leq f_2$ on $E$; monotonicity gives $\int_Ef^\pm\,d\mu<+\infty$.
\item Since $f^+-f^-=f_1-f_2$ on $E$, $f^++f_2=f_1+f^-$ on $E$.\\
Nonnegative additivity gives $\int_E f^+\,d\mu+\int_E f_2\,d\mu=\int_E f_1\,d\mu+\int_E f^-\,d\mu$; subtract the finite terms.
\end{enumerate}
\hlt{Structural Commentary}: Finite positive decompositions determine one signed integral.\\
\hlt{Key Insight}: Move all nonnegative terms to one equality, integrate, then cancel finite terms.
\end{proof}
\end{proposition}

Simple functions are finite linear combinations of indicators of measurable sets: coefficients carry size, while the level sets carry measurability.

\begin{lemma} \label{lem:simple-integral-refinement} \hlt{Refining a Simple Integral}\\
Let $s\geq0$ be simple measurable and suppose $s=\sum\limits_{k=1}^{m}b_k\mathbf{1}_{F_k}$ with $b_k\geq0$ and $F_k\in\mathcal{M}$.\\
Then $\forall E\in\mathcal{M}$, $\mathcal{I}_E(s)=\sum\limits_{k=1}^{m}b_k\mu(E\cap F_k)$.
\begin{proof}
For $\eta\in\{0,1\}^m$, set $F_k^1\defeq F_k$, $F_k^0\defeq F_k^c$, discard empty atoms $G_\eta\defeq\bigcap\limits_{k=1}^{m}F_k^{\eta_k}$. On $G_\eta$, $s=c_\eta\defeq\sum\limits_{\eta_k=1}b_k$.\\
Group equal $c_\eta$ values to recover the standard representation; finite additivity gives\\
$\mathcal{I}_E(s)=\sum\limits_\eta c_\eta\mu(E\cap G_\eta) =\sum\limits_{k=1}^{m}b_k\sum\limits_{\eta:\eta_k=1}\mu(E\cap G_\eta) =\sum\limits_{k=1}^{m}b_k\mu(E\cap F_k)$.\\
\hlt{Structural Commentary}: Refine nonstandard sums to one finite partition.\\
\hlt{Key Insight}: Refine first, then regroup by equal function values.
\end{proof}
\end{lemma}

\begin{proposition} \label{prop:simple-integral-linearity-and-measure} \hlt{Simple Integral Algebra}\\
Let $(X,\mathcal{M},\mu)$ be a measure space, let $E\in\mathcal{M}$, let $c\geq0$, and let $s,t\in\mathcal{S}^+$. For fixed $s\in\mathcal{S}^+$, define $\nu_s:\mathcal{M}\to[0,+\infty]$ by $\nu_s(A)\defeq\mathcal{I}_A(s)$. Then:
\begin{enumerate}[label=(\roman*)]
\setlength{\itemsep}{0pt}
\item $\mathcal{I}_E(cs)=c\mathcal{I}_E(s)$.
\item $\mathcal{I}_E(s+t)=\mathcal{I}_E(s)+\mathcal{I}_E(t)$.
\item $\nu_s$ is a measure on $\mathcal{M}$.
\end{enumerate}
\begin{proof} Write $s=\sum\limits_ia_i\mathbf{1}_{A_i}$ and $t=\sum\limits_jb_j\mathbf{1}_{B_j}$ on finite measurable partitions of $X$, with $a_i,b_j\geq0$.
\begin{enumerate}[label=(\roman*)]
\setlength{\itemsep}{0pt}
\item $\mathcal{I}_E(cs)=\sum\limits_i(ca_i)\mu(E\cap A_i)=c\sum\limits_ia_i\mu(E\cap A_i)=c\mathcal{I}_E(s)$, with $0\cdot(+\infty)=0$.
\item On the common refinement $A_i\cap B_j$, $s+t=\sum\limits_{i,j}(a_i+b_j)\mathbf{1}_{A_i\cap B_j}$.\\
$\mathcal{I}_E(s+t)=\sum\limits_{i,j}(a_i+b_j)\mu(E\cap A_i\cap B_j)=\sum\limits_ia_i\mu(E\cap A_i)+\sum\limits_jb_j\mu(E\cap B_j)$ by Lemma \ref{lem:simple-integral-refinement}.
\item For $A\in\mathcal{M}$, $\nu_s(A)=\sum\limits_ia_i\mu(A\cap A_i)$.\\
Each $A\mapsto\mu(A\cap A_i)$ is a measure, and finite nonnegative sums of measures are measures.
\end{enumerate}
\hlt{Structural Commentary}: Simple integration is weighted measure summation on finite partitions.\\
\hlt{Key Insight}: Refine to common atoms, then use finite additivity.
\end{proof}
\end{proposition}

\begin{proposition} \label{prop:positive-integral-monotonicity} \hlt{Monotonicity for Nonnegative Integrals}\\
Let $f,g\geq0$ be measurable and let $E,F\in\mathcal{M}$.
\begin{enumerate}[label=(\roman*)]
\setlength{\itemsep}{0pt}
\item If $f\leq g$ on $E$, then $\int_E f\,d\mu\leq\int_E g\,d\mu$.
\item If $E\subseteq F$, then $\int_E f\,d\mu\leq\int_F f\,d\mu$.
\end{enumerate}
\begin{proof}
\begin{enumerate}[label=(\roman*)]
\setlength{\itemsep}{0pt}
\item If $s\in\mathcal{S}^+(f)$, then $s\mathbf{1}_E\in\mathcal{S}^+(g)$, $\mathcal{I}_E(s)=\mathcal{I}_E(s\mathbf{1}_E)\leq\int_Eg\,d\mu$. Taking $\sup\limits_{s\in\mathcal{S}^+(f)}$ gives $\int_Ef\,d\mu\leq\int_Eg\,d\mu$.
\item If $E\subseteq F$ and $s\in\mathcal{S}^+(f)$, then $\mathcal{I}_E(s)\leq\mathcal{I}_F(s)$. Taking $\sup\limits_{s\in\mathcal{S}^+(f)}$ gives $\int_Ef\,d\mu\leq\int_Ff\,d\mu$.
\end{enumerate}
\hlt{Structural Commentary}: The nonnegative integral is a supremum over simple functions below the integrand.\\
\hlt{Key Insight}: More room below the integrand, or over a larger set, can only increase the supremum.
\end{proof}
\end{proposition}

\begin{proposition} \label{prop:chebyshev-inequality-nonnegative-integral} \hlt{Chebyshev Inequality for Nonnegative Integrals}\\
Let $(X,\mathcal{M},\mu)$ be a measure space, let $E\in\mathcal{M}$, and let $f:X\to[0,+\infty]$ be measurable.\\
For $\lambda>0$, set $E_\lambda\defeq E\cap\{f\geq\lambda\}$. Then $E_\lambda\in\mathcal{M}$ and $\mu(E_\lambda)\leq\lambda^{-1}\int_Ef\,d\mu$.
\begin{proof} Measurability gives $E_\lambda\in\mathcal{M}$.\\
Since $\lambda\mathbf{1}_{E_\lambda}\leq f$ on $E$, monotonicity gives $\lambda\mu(E_\lambda)=\int_E\lambda\mathbf{1}_{E_\lambda}\,d\mu\leq\int_Ef\,d\mu$; divide by $\lambda>0$.\\
\hlt{Structural Commentary}: Level-set minorants convert mass to measure.\\
\hlt{Key Insight}: Use $\lambda\mathbf{1}_{E_\lambda}\leq f$ on $E$.
\end{proof}
\end{proposition}

\begin{method} \label{meth:splitting-simple-minorants-under-sum} \hlt{Splitting a Simple Minorant Under a Sum}\\
Let $E\in\mathcal{M}$, $f,g\geq0$ measurable, and $\omega\in\mathcal{S}^+(f+g)$ simple.
\begin{enumerate}[label=(\roman*)]
\setlength{\itemsep}{0pt}
\item \hlt{Grid}: for $n\geq1$, use $D_n\defeq\{0,\frac{1}{n},\ldots,1\}$.
\item \hlt{Split}: define $\varphi_n(x)\defeq\max\{\theta\omega(x):\theta\in D_n,\ \theta\omega(x)\leq f(x)\}$; $\psi_n(x)\defeq\left(\left(1-\frac{1}{n}\right)\omega(x)-\varphi_n(x)\right)^+$.
\item \hlt{Bounds}: $\varphi_n,\psi_n$ are simple, $\varphi_n\leq f$, $\psi_n\leq g$, and $(1-\frac{1}{n})\omega\leq\varphi_n+\psi_n$.
\item \hlt{Use}: simple additivity gives $(1-\frac{1}{n})\mathcal{I}_E(\omega)\leq\int_Ef\,d\mu+\int_Eg\,d\mu$; then let $n\to\infty$ and take $\sup\limits_{\omega\in\mathcal{S}^+(f+g)}$.
\end{enumerate}
\hlt{Key Insight}: Grid-split $\omega$ between $f$ and $g$, losing only $\frac{1}{n}$.
\end{method}

\begin{proposition} \label{prop:nonnegative-integral-homogeneity-additivity} \hlt{Homogeneity and Additivity for Nonnegative Integrals}\\
Let $f,g\geq0$ be measurable, let $c\geq0$, and let $E\in\mathcal{M}$. All identities are in $[0,+\infty]$, with $0\cdot(+\infty)=0$.
\begin{enumerate}[label=(\roman*)]
\setlength{\itemsep}{0pt}
\item \hlt{Homogeneity}: $\int_E cf\,d\mu=c\int_E f\,d\mu$.
\item \hlt{Additivity}: $\int_E(f+g)\,d\mu=\int_Ef\,d\mu+\int_Eg\,d\mu$.
\end{enumerate}
\begin{proof}
\begin{enumerate}[label=(\roman*)]
\setlength{\itemsep}{0pt}
\item If $c=0$, both sides are $0$; if $c>0$, $s\mapsto cs$ bijects $\mathcal{S}^+(f)$ with $\mathcal{S}^+(cf)$ and satisfies $\mathcal{I}_E(cs)=c\mathcal{I}_E(s)$.\\
Hence taking suprema gives the claim.
\item If $s\in\mathcal{S}^+(f)$ and $t\in\mathcal{S}^+(g)$, then $s+t\in\mathcal{S}^+(f+g)$ and $\mathcal{I}_E(s)+\mathcal{I}_E(t)=\mathcal{I}_E(s+t)\leq\int_E(f+g)\,d\mu$, so taking $\sup\limits_{s,t}$ gives $\int_Ef\,d\mu+\int_Eg\,d\mu\leq\int_E(f+g)\,d\mu$.\\
Conversely, for $\omega\in\mathcal{S}^+(f+g)$, choose simple $\varphi_n\leq f$ and $\psi_n\leq g$ with $(1-\frac{1}{n})\omega\leq\varphi_n+\psi_n$.\\
Then $\left(1-\frac{1}{n}\right)\mathcal{I}_E(\omega)\leq\mathcal{I}_E(\varphi_n)+\mathcal{I}_E(\psi_n)\leq\int_E f\,d\mu+\int_E g\,d\mu$. Let $n\to\infty$, then take $\sup\limits_{\omega\in\mathcal{S}^+(f+g)}$.
\end{enumerate}
\hlt{Structural Commentary}: The nonnegative integral is linear before signed integrability is introduced.\\
\hlt{Key Insight}: Scale minorants; split sum-minorants on a finite grid.
\end{proof}
\end{proposition}

\begin{perspective} \label{per:riemann-step-vs-lebesgue-simple} \hlt{Step Functions versus Simple Functions}\\
Riemann: $\sup\sum\limits_i m_i\ell(I_i)$ over interval partitions. Lebesgue: $\sup\{\mathcal{I}(s):0\leq s\leq f,\ s=\sum\limits_ic_i\mathbf{1}_{E_i},\ E_i\in\mathcal{M}\}$.\\
\hlt{Structural Commentary}: Riemann partitions intervals; Lebesgue partitions level sets.\\
\hlt{Key Insight}: Same finite constants, measurable pieces.
\end{perspective}

\begin{perspective} \label{per:undergraph-geometric-integral} \hlt{Undergraph Picture for Nonnegative Functions}\\
For measurable $f:\R\to[0,+\infty]$, set $\mathcal{U}_f\defeq\{(x,y)\in\R\times[0,+\infty):0\leq y<f(x)\}$.\\
Geometrically, $\int f\,dm=m_2(\mathcal{U}_f)$; analytically, the integral is built as $\sup\mathcal{I}(s)$ over simple $0\leq s\leq f$.
\end{perspective}

\begin{figure}[H]
  \centering
  \begin{tikzpicture}[
  scale=0.82,
  >=Stealth,
  every node/.style={font=\scriptsize},
  axis/.style={->, black!40, line width=0.5pt},
  curve/.style={blue!70!black, line width=0.85pt},
  under/.style={fill=blue!18, draw=none}
]
  \draw[axis] (-0.35,0) -- (7.75,0) node[right] {$x$};
  \draw[axis] (0,-0.18) -- (0,3.25) node[above] {$y$};

  \path[under]
    (0,0)
    -- plot[domain=0:7.20, samples=120]
      (\x,{1.50 + 0.30*sin(105*\x) + 0.40*sin(51*\x) + 0.18*cos(180*\x)})
    -- (7.20,0) -- cycle;

  \draw[curve, domain=0:7.20, samples=120]
      plot (\x,{1.50 + 0.30*sin(105*\x) + 0.40*sin(51*\x) + 0.18*cos(180*\x)});

  \node[blue!45!black] at (3.85,0.55) {$\mathcal{U}f$};
  \node[blue!70!black, fill=white, inner sep=1pt] at (6.65,2.02) {$f$};
\end{tikzpicture}

  \caption{The integral of $f\geq0$ is modeled by the measure of its undergraph.}
  \label{fig:lebesgue-undergraph}
\end{figure}

\begin{proposition} \label{prop:basic-lebesgue-integral-properties} \hlt{Basic Integral Properties}\\
Let $E\in\mathcal{M}$.
\begin{enumerate}[label=(\roman*)]
\setlength{\itemsep}{0pt}
\item If $f$ is measurable and bounded on $E$, and $\mu(E)<+\infty$, then $f\in\mathscr{L}^1(\mu)$ on $E$.
\item If $a\leq f(x)\leq b$ on $E$ and $\mu(E)<+\infty$, then $a\mu(E)\leq\int_E f\,d\mu\leq b\mu(E)$.
\item If $f,g\in\mathscr{L}^1(\mu)$ on $E$ and $f\leq g$ on $E$, then $\int_E f\,d\mu\leq\int_E g\,d\mu$.
\item If $f\in\mathscr{L}^1(\mu)$ on $E$, then $cf\in\mathscr{L}^1(\mu)$ on $E$ for finite $c$, and $\int_E cf\,d\mu=c\int_E f\,d\mu$.
\item If $\mu(E)=0$ and $f$ is measurable, then $\int_E f\,d\mu=0$.
\item If $f\in\mathscr{L}^1(\mu)$ on $E$, $A\in\mathcal{M}$, and $A\subseteq E$, then $f\in\mathscr{L}^1(\mu)$ on $A$.
\end{enumerate}
\begin{proof}
\begin{enumerate}[label=(\roman*)]
\setlength{\itemsep}{0pt}
\item If $\abs{f}\leq M$ on $E$, then $0\leq f^\pm\leq M\mathbf{1}_E$ on $E$, so $\int_Ef^\pm\,d\mu\leq M\mu(E)<+\infty$.
\item Since $a\mathbf{1}_E\leq f\leq b\mathbf{1}_E$ on $E$, order for the positive integral gives $a\mu(E)\leq\int_Ef\,d\mu\leq b\mu(E)$.
\item Since $g-f\geq0$ on $E$, $\int_E(g-f)\,d\mu\geq0$, hence $\int_Ef\,d\mu\leq\int_Eg\,d\mu$.
\item For $c\geq0$, scale the simple approximants. For $c<0$, use $(cf)^+=(-c)f^-$ and $(cf)^-=(-c)f^+$.
\item If $\mu(E)=0$, then $\forall s\geq0$ simple, $\mathcal{I}_E(s)=0$, so $\int_Ef^\pm\,d\mu=0$.
\item Since $A\subseteq E$, $\mathcal{I}_A(s)\leq \mathcal{I}_E(s)$ for $s\geq0$ simple; hence $\int_Af^\pm\,d\mu\leq\int_Ef^\pm\,d\mu<+\infty$.
\end{enumerate}
\hlt{Structural Commentary}: Estimate $f^\pm$ by simpler functions.\\
\hlt{Key Insight}: Monotonicity makes bounds, null sets, and restrictions immediate.
\end{proof}
\end{proposition}

\begin{theorem} \label{thm:integral-countably-additive-set-function} \hlt{The Integral as a Measure or Charge}\\
For $A\in\mathcal{M}$, define $\phi_f(A)\defeq\int_Af\,d\mu$ when $f\geq0$ is measurable, and define $\lambda_f(A)\defeq\int_Af\,d\mu$ when $f\in\mathscr{L}^1(\mu)$.
\begin{enumerate}[label=(\roman*)]
\setlength{\itemsep}{0pt}
\item If $f\geq0$ is measurable, then $\phi_f:\mathcal{M}\to[0,+\infty]$ is a measure.
\item If $f\in\mathscr{L}^1(\mu)$, then $\lambda_f:\mathcal{M}\to\R$ is finite-valued and countably additive.
\end{enumerate}
\begin{proof}
\begin{enumerate}[label=(\roman*)]
\setlength{\itemsep}{0pt}
\item For simple $s\geq0$, $A\mapsto\mathcal{I}_A(s)$ is a measure, so $\phi_f(\emptyset)=0$.\\
If $A=\bigsqcup\limits_nA_n$, $s\in\mathcal{S}^+(f)$, then $\mathcal{I}_A(s)=\sum\limits_n\mathcal{I}_{A_n}(s)\leq\sum\limits_n\phi_f(A_n)$; taking $\sup\limits_s$ gives $\phi_f(A)\leq\sum\limits_n\phi_f(A_n)$. Conversely, $\forall N$, finite additivity gives $\sum\limits_{n=1}^{N}\phi_f(A_n)=\phi_f(\bigcup\limits_{n=1}^{N}A_n)\leq\phi_f(A)$.\\
Let $N\to\infty$. Hence $\phi_f(A)=\sum\limits_n\phi_f(A_n)$.
\item Since $f\in\mathscr{L}^1(\mu)$, $\phi_{f^+}(X),\phi_{f^-}(X)<+\infty$ and $\lambda_f=\phi_{f^+}-\phi_{f^-}$ is finite-valued.\\
If $A=\bigsqcup\limits_nA_n$, then item (i) gives $\lambda_f(A)=\sum\limits_n\phi_{f^+}(A_n)-\sum\limits_n\phi_{f^-}(A_n)=\sum\limits_n\lambda_f(A_n)$.
\end{enumerate}
\hlt{Structural Commentary}: Setwise integration turns densities into measures.\\
\hlt{Key Insight}: Sup simple measures; subtract finite signs.
\end{proof}
\end{theorem}

\begin{corollary} \label{cor:integration-countable-additivity-and-continuity-over-domains} \hlt{Domain Additivity and Continuity for Integrable Functions}\\
Let $(X,\mathcal{M},\mu)$ be a measure space, let $E\in\mathcal{M}$, let $f\in\mathscr{L}^1(\mu)$ on $E$, and set $\mathcal{M}_E\defeq\{A\in\mathcal{M}:A\subseteq E\}$.
\begin{enumerate}[label=(\roman*)]
\setlength{\itemsep}{0pt}
\item If $(A_n)\subseteq\mathcal{M}_E$ are pairwise disjoint, then $\int_{\bigsqcup\limits_nA_n}f\,d\mu=\sum\limits_n\int_{A_n}f\,d\mu$ and $\sum\limits_n\abs{\int_{A_n}f\,d\mu}<+\infty$.
\item If $A_n,A\in\mathcal{M}_E$ and $A_n\uparrow A$, then $\int_Af\,d\mu=\lim\limits_{n\to\infty}\int_{A_n}f\,d\mu$.
\item If $A_n,A\in\mathcal{M}_E$ and $A_n\downarrow A$, then $\int_Af\,d\mu=\lim\limits_{n\to\infty}\int_{A_n}f\,d\mu$.
\end{enumerate}
\begin{proof} Let $\mu_E\defeq\mu|_{\mathcal{M}_E}$ and define $\lambda(A)\defeq\int_Af\,d\mu$, $\nu(A)\defeq\int_A\abs f\,d\mu$ on $\mathcal{M}_E$.\\
By Theorem \ref{thm:integral-countably-additive-set-function}, $\lambda$ is a finite signed measure, $\nu(E)<+\infty$, and $\abs{\lambda(A)}\leq\nu(A)$.
\begin{enumerate}[label=(\roman*)]
\setlength{\itemsep}{0pt}
\item Additivity gives $\lambda(\bigsqcup\limits_nA_n)=\sum\limits_n\lambda(A_n)$, and $\sum\limits_n\abs{\lambda(A_n)}\leq\sum\limits_n\nu(A_n)=\nu(\bigsqcup\limits_nA_n)\leq\nu(E)<+\infty$.
\item Set $B_1=A_1$ and $B_n=A_n\setminus A_{n-1}$ for $n\geq2$.\\
Since $A=\bigsqcup\limits_nB_n$ and $A_N=\bigsqcup\limits_{n\leq N}B_n$, (i) gives $\lambda(A_N)=\sum\limits_{n\leq N}\lambda(B_n)\to\sum\limits_n\lambda(B_n)=\lambda(A)$.
\item Set $C_n\defeq A_1\setminus A_n$ and $C\defeq A_1\setminus A$. Then $C_n\uparrow C$, so (ii) gives $\lambda(C_n)\to\lambda(C)$.\\
Hence $\lambda(A_n)=\lambda(A_1)-\lambda(C_n)\to\lambda(A_1)-\lambda(C)=\lambda(A)$.
\end{enumerate}
\hlt{Structural Commentary}: Domain integrals form a finite signed measure.\\
\hlt{Key Insight}: Monotone domain limits are signed-measure continuity.
\end{proof}
\end{corollary}

\begin{proposition} \label{prop:finite-measure-support-integrals} \hlt{Finite-Measure Support}\\
Let $(X,\mathcal{M},\mu)$ be a measure space, let $E\in\mathcal{M}$, and let $h:E\to\R$ be bounded and measurable.\\
Suppose $h=0$ on $E\setminus E_0$ for some $E_0\in\mathcal{M}$ with $E_0\subseteq E$ and $\mu(E_0)<+\infty$. Then:
\begin{enumerate}[label=(\roman*)]
\setlength{\itemsep}{0pt}
\item $h\in\mathscr{L}^1(\mu)$ on $E$.
\item $\int_Eh\,d\mu=\int_{E_0}h\,d\mu$.
\item If $E_1\in\mathcal{M}$, $E_1\subseteq E$, $\mu(E_1)<+\infty$, and $h=0$ on $E\setminus E_1$, then $\int_{E_0}h\,d\mu=\int_{E_1}h\,d\mu$.
\end{enumerate}
\begin{proof}
\begin{enumerate}[label=(\roman*)]
\setlength{\itemsep}{0pt}
\item Choose $M<+\infty$ with $\abs h\leq M$ on $E$.\\
Since $h=0$ on $E\setminus E_0$, $0\leq h^\pm\leq M\mathbf{1}_{E_0}$ on $E$, so $\int_Eh^\pm\,d\mu\leq M\mu(E_0)<+\infty$.
\item Since $\int_{E\setminus E_0}h\,d\mu=0$, $\int_Eh\,d\mu=\int_{E_0}h\,d\mu$ by Corollary \ref{cor:integration-countable-additivity-and-continuity-over-domains}.
\item Apply item (ii) with $E_0$ and with $E_1$: $\int_{E_0}h\,d\mu=\int_Eh\,d\mu=\int_{E_1}h\,d\mu$.
\end{enumerate}
\hlt{Structural Commentary}: Finite support turns boundedness into integrability.\\
\hlt{Key Insight}: The integral only sees where $h$ can be nonzero.
\end{proof}
\end{proposition}

\begin{corollary} \label{cor:compact-continuous-finite-measure-integrable} \hlt{Compact Continuous Functions on Finite Measure Spaces}\\
Let $X$ be a compact topological space, let $\mathcal{M}$ be a $\sigma$-algebra containing the topology of $X$, and let $\mu(X)<+\infty$.\\
If $f:X\to\R$ is continuous, then:
\begin{enumerate}[label=(\roman*)]
\setlength{\itemsep}{0pt}
\item $f$ is measurable.
\item $f\in\mathscr{L}^1(\mu)$.
\end{enumerate}
\begin{proof}
\begin{enumerate}[label=(\roman*)]
\setlength{\itemsep}{0pt}
\item $\forall O\subseteq\R$ open, $f^{-1}(O)$ is open in $X$ by continuity, hence $f^{-1}(O)\in\mathcal{M}$.
\item The continuous image $f(X)$ is compact in $\R$, hence bounded.\\
Since $\mu(X)<+\infty$, $f\in\mathscr{L}^1(\mu)$ by Proposition \ref{prop:basic-lebesgue-integral-properties}(i).
\end{enumerate}
\hlt{Structural Commentary}: Compactness supplies boundedness, while the measure hypothesis supplies finite mass.\\
\hlt{Key Insight}: Continuous on compact plus finite measure implies bounded measurable, hence integrable.
\end{proof}
\end{corollary}

\begin{corollary} \label{cor:integral-null-and-finite-ae-consequences} \hlt{Null and Finiteness Consequences for $f\geq0$}\\
Let $f:X\to[0,+\infty]$ be measurable and $\phi_f(A)\defeq\int_A f\,d\mu$.
\begin{enumerate}[label=(\roman*)]
\setlength{\itemsep}{0pt}
\item $\mu(A)=0\Longrightarrow\phi_f(A)=0$ for $A\in\mathcal{M}$.
\item $\int_X f\,d\mu=0\Longleftrightarrow f=0$ a.e.
\item $\int_X f\,d\mu<+\infty\Longrightarrow\mu\{f=+\infty\}=0$ and $\{f>0\}$ is $\sigma$-finite.
\end{enumerate}
\begin{proof}
\begin{enumerate}[label=(\roman*)]
\setlength{\itemsep}{0pt}
\item[$(i)$] If $\mu(A)=0$ and $0\leq s\leq f$ is simple, then $\mathcal{I}_A(s)=0$; hence $\phi_f(A)=\sup\limits_s\mathcal{I}_A(s)=0$.
\item[$(ii)$] ($\Rightarrow$) Set $E_n\defeq\{f>\frac{1}{n}\}$. By Proposition \ref{prop:chebyshev-inequality-nonnegative-integral}, $\mu(E_n)\leq n\int_X f\,d\mu=0$.\\
Thus $\forall n,\ \mu(E_n)=0$, and $\{f>0\}=\bigcup\limits_nE_n$ is null.
\item[$(ii)$] ($\Leftarrow$) If $N\defeq\{f>0\}$ null, every simple $0\leq s\leq f$ vanishes on $X\setminus N$, so $\mathcal{I}_X(s)=\mathcal{I}_N(s)=0$; take $\sup\limits_s$.
\item[$(iii)$] Let $F_n\defeq\{f\geq n\}$. By Proposition \ref{prop:chebyshev-inequality-nonnegative-integral}, $\mu(F_n)\leq n^{-1}\int_X f\,d\mu$, so $\mu\{f=+\infty\}\leq\mu(F_n)\to0$.\\
Also $\{f>0\}=\bigcup\limits_n\{f>\frac{1}{n}\}$ and $\mu\{f>\frac{1}{n}\}\leq n\int_X f\,d\mu<+\infty$.
\end{enumerate}
\hlt{Structural Commentary}: Finite nonnegative integrals control threshold sets.\\
\hlt{Key Insight}: Threshold sets convert integral bounds into measure bounds.
\end{proof}
\end{corollary}

\begin{theorem} \label{thm:absolute-integrability-and-domination} \hlt{Absolute Integrability and Domination}\\
Let $(X,\mathcal{M},\mu)$ be a measure space and let $E\in\mathcal{M}$. Then:
\begin{enumerate}[label=(\roman*)]
\setlength{\itemsep}{0pt}
\item If $f\in\mathscr{L}^1(\mu)$ on $E$, then $\abs{f}\in\mathscr{L}^1(\mu)$ on $E$ and $\abs{\int_E f\,d\mu}\leq\int_E\abs{f}\,d\mu$.
\item If $f:X\to\overline{\R}$ is measurable, $g\in\mathscr{L}^1(\mu)$ on $E$, and $\abs{f}\leq g$ on $E$, then $f\in\mathscr{L}^1(\mu)$ on $E$.
\end{enumerate}
\begin{proof}
\begin{enumerate}[label=(\roman*)]
\setlength{\itemsep}{0pt}
\item Set $a\defeq\int_Ef^+\,d\mu$ and $b\defeq\int_Ef^-\,d\mu$. Since $f\in\mathscr{L}^1(\mu)$ on $E$, $a,b<+\infty$.\\
Hence $\int_E\abs f\,d\mu=\int_E(f^++f^-)\,d\mu=a+b<+\infty$, and $\abs{\int_Ef\,d\mu}=\abs{a-b}\leq a+b=\int_E\abs f\,d\mu$.
\item Since $\abs f\leq g$ on $E$, $g\geq0$ on $E$ and $\int_Eg\,d\mu<+\infty$.\\
Also $0\leq f^\pm\leq\abs f\leq g$ on $E$, so monotonicity gives $\int_Ef^\pm\,d\mu<+\infty$.
\end{enumerate}
\hlt{Structural Commentary}: Both claims reduce to finite positive and negative parts.\\
\hlt{Key Insight}: Use $\abs f=f^++f^-$ and $0\leq f^\pm\leq g$.
\end{proof}
\end{theorem}

\begin{proposition} \label{prop:lebesgue-integral-linearity} \hlt{Linearity of the Lebesgue Integral}\\
Let $(X,\mathcal{M},\mu)$ be a measure space, let $E\in\mathcal{M}$, let $f,g\in\mathscr{L}^1(\mu)$ on $E$, and let $a,b\in\R$.\\
Then:
\begin{enumerate}[label=(\roman*)]
\setlength{\itemsep}{0pt}
\item For $c\in\R$ and $u\in\mathscr{L}^1(\mu)$ on $E$, $cu\in\mathscr{L}^1(\mu)$ on $E$ and $\int_Ecu\,d\mu=c\int_Eu\,d\mu$.
\item $f+g\in\mathscr{L}^1(\mu)$ on $E$ and $\int_E(f+g)\,d\mu=\int_Ef\,d\mu+\int_Eg\,d\mu$.
\item $af+bg\in\mathscr{L}^1(\mu)$ on $E$ and $\int_E(af+bg)\,d\mu=a\int_Ef\,d\mu+b\int_Eg\,d\mu$.
\end{enumerate}
\begin{proof}
\begin{enumerate}[label=(\roman*)]
\setlength{\itemsep}{0pt}
\item Since $(cu)^\pm=c^+u^\pm+c^-u^\mp$, both $(cu)^+$ and $(cu)^-$ have finite integral. Nonnegative linearity gives $\int_Ecu\,d\mu=(c^+a_u+c^-b_u)-(c^+b_u+c^-a_u)=c(a_u-b_u)$, where $a_u\defeq\int_Eu^+\,d\mu$ and $b_u\defeq\int_Eu^-\,d\mu$.
\item Since $\abs{f+g}\leq\abs f+\abs g$ and $\abs f+\abs g\in\mathscr{L}^1(\mu)$ on $E$, $f+g\in\mathscr{L}^1(\mu)$ on $E$ by Theorem \ref{thm:absolute-integrability-and-domination}(ii).\\
Integrate $(f+g)^+ + f^- + g^- = (f+g)^- + f^+ + g^+$ over $E$ and cancel finite terms.
\item By (i), $af,bg\in\mathscr{L}^1(\mu)$ on $E$, $\int_Eaf\,d\mu=a\int_Ef\,d\mu$, and $\int_Ebg\,d\mu=b\int_Eg\,d\mu$.\\
Applying (ii) to $af,bg$ gives $af+bg\in\mathscr{L}^1(\mu)$ on $E$ and $\int_E(af+bg)\,d\mu=a\int_Ef\,d\mu+b\int_Eg\,d\mu$.
\end{enumerate}
\hlt{Structural Commentary}: Signed linearity follows from finite sign pieces.\\
\hlt{Key Insight}: Integrate one nonnegative equality, then cancel finite terms.
\end{proof}
\end{proposition}

\begin{proposition} \label{prop:l1-level-sets-null-and-indefinite-integrals} \hlt{$L^1$ Level Sets, Null Changes, and Indefinite Integrals}\\
Let $f,g\in\mathscr{L}^1(\mu)$ on $X$, and write $\lambda_u(A)\defeq\int_Au\,d\mu$.
\begin{enumerate}[label=(\roman*)]
\setlength{\itemsep}{0pt}
\item For $a>0$, $\mu\{\abs{f}\geq a\}\leq a^{-1}\int_X\abs{f}\,d\mu<+\infty$; hence $\{f\neq0\}$ is $\sigma$-finite.
\item If measurable $h=0$ a.e., then $h\in\mathscr{L}^1(\mu)$ and $\lambda_h\equiv0$.
\item $\lambda_f\geq0$ on $\mathcal{M}\Longrightarrow f\geq0$ a.e.; $\lambda_f\equiv0\Longrightarrow f=0$ a.e.
\item $\lambda_f=\lambda_g$ on $\mathcal{M}\Longleftrightarrow f=g$ a.e.
\end{enumerate}
\begin{proof}
\begin{enumerate}[label=(\roman*)]
\setlength{\itemsep}{0pt}
\item Since $a\mathbf{1}_{\{\abs{f}\geq a\}}\leq\abs{f}$, integration gives $a\mu\{\abs{f}\geq a\}\leq\int_X\abs{f}\,d\mu$. Also $\{f\neq0\}=\bigcup\limits_{k\geq1}\{\abs{f}\geq\frac{1}{k}\}$.
\item Let $N\defeq\{h\neq0\}$. For $A\in\mathcal{M}$, $h^\pm\mathbf{1}_A$ vanish off $A\cap N$, a null set.\\
Hence $\int_Ah^+\,d\mu=\int_Ah^-\,d\mu=0$, so $\lambda_h(A)=0$.
\item For $k\geq1$, set $N_k\defeq\{f\leq-\frac{1}{k}\}$. By (i), $\mu(N_k)<+\infty$, and $\lambda_f(N_k)\leq-k^{-1}\mu(N_k)$.\\
If $\lambda_f\geq0$, then $\mu(N_k)=0$, so $\{f<0\}=\bigcup\limits_kN_k$ is null. If $\lambda_f\equiv0$, apply this to $f$ and $-f$.
\item If $\lambda_f=\lambda_g$, then $\lambda_{f-g}\equiv0$, so $f-g=0$ a.e. by (iii). Conversely, if $f=g$ a.e., apply (ii) to $h=f-g$.
\end{enumerate}
\hlt{Structural Commentary}: An $L^1$ function is detected by its integrals over all measurable sets, up to null sets.\\
\hlt{Key Insight}: Test sign on the level sets where the function is bounded away from $0$.
\end{proof}
\end{proposition}

\begin{proposition} \label{prop:l1-simple-approximation-and-bounded-multipliers} \hlt{$L^1$ Approximation and Bounded Multipliers}\\
Let $f\in\mathscr{L}^1(\mu)$ on $X$.
\begin{enumerate}[label=(\roman*)]
\setlength{\itemsep}{0pt}
\item $\forall\varepsilon>0\ \exists$ a measurable simple function $\varphi$ with $\int_X\abs{f-\varphi}\,d\mu<\varepsilon$.
\item If $g$ is bounded measurable, then $fg\in\mathscr{L}^1(\mu)$.
\end{enumerate}
\begin{proof}
\begin{enumerate}[label=(\roman*)]
\setlength{\itemsep}{0pt}
\item Choose simple $0\leq s_n\uparrow f^+$ and $0\leq t_n\uparrow f^-$. Then $\int_X(f^+-s_n)\,d\mu\to0$ and $\int_X(f^--t_n)\,d\mu\to0$.\\
For large $n$, $\varphi\defeq s_n-t_n$ satisfies $\int_X\abs{f-\varphi}\,d\mu\leq\int_X(f^+-s_n)\,d\mu+\int_X(f^--t_n)\,d\mu<\varepsilon$.
\item If $\abs{g}\leq M$, then $fg$ is measurable and $\int_X\abs{fg}\,d\mu\leq M\int_X\abs{f}\,d\mu<+\infty$.
\end{enumerate}
\hlt{Structural Commentary}: $L^1$ is simple-closed and bounded-multiplier stable.\\
\hlt{Key Insight}: Approximate signs separately; dominate products by $M\abs{f}$.
\end{proof}
\end{proposition}

\begin{proposition} \label{prop:finite-measure-indicator-continuous-l1-approximation} \hlt{Finite-Measure Indicators by Intervals and Continuous Functions}\\
Let $A\subseteq\R$ be Lebesgue measurable with $m(A)<+\infty$. Then $\forall\varepsilon>0$:
\begin{enumerate}[label=(\roman*)]
\setlength{\itemsep}{0pt}
\item $\exists G=\bigcup\limits_{j=1}^{N}I_j$, $I_j$ open intervals, such that $m(A\triangle G)<\varepsilon$.
\item $\exists f\in C(\R)$ such that $0\leq f\leq1$ and $\int_{\R}\abs{\mathbf{1}_A-f}\,dm<\varepsilon$.
\end{enumerate}
\begin{proof}
\begin{enumerate}[label=(\roman*)]
\setlength{\itemsep}{0pt}
\item Choose open $O\supseteq A$ with $m(O\setminus A)<\frac{\varepsilon}{2}$ and $m(O)<+\infty$.\\
Write $O=\bigsqcup\limits_{j\geq1}I_j$ with open intervals $I_j$, choose $N$ with $m(O\setminus G)<\frac{\varepsilon}{2}$ for $G\defeq\bigcup\limits_{j=1}^{N}I_j$.\\
Then $A\setminus G\subseteq O\setminus G$ and $G\setminus A\subseteq O\setminus A$, so $m(A\triangle G)<\varepsilon$.
\item Apply (i) with $\frac{\varepsilon}{2}$ to get finite open $G$. Choose compact $K\subseteq G$ with $m(G\setminus K)<\frac{\varepsilon}{2}$, and choose piecewise-linear $f$ with $0\leq f\leq1$, $f=1$ on $K$, $f=0$ on $\R\setminus G$.\\
Then $\int_{\R}\abs{\mathbf{1}_G-f}\,dm\leq m(G\setminus K)<\frac{\varepsilon}{2}$, hence $\int_{\R}\abs{\mathbf{1}_A-f}\,dm\leq m(A\triangle G)+\int_{\R}\abs{\mathbf{1}_G-f}\,dm<\varepsilon$.
\end{enumerate}
\hlt{Structural Commentary}: Finite-measure sets can be made geometrically simple up to small $L^1$ error.\\
\hlt{Key Insight}: Approximate by intervals, then smooth the indicator.
\end{proof}
\end{proposition}

\begin{example} \label{ex:counting-and-unit-mass-integrals} \hlt{Counting and Unit-Mass Integrals}
\begin{enumerate}[label=(\roman*)]
\setlength{\itemsep}{0pt}
\item \hlt{Counting}: for counting measure $\#$ on $(\N,2^\N)$ and $f:\N\to\R$,\\
$f\in\mathscr{L}^1(\#)\Longleftrightarrow\sum\limits_{n=1}^{\infty}\abs{f(n)}<+\infty$, and then $\int_{\N}f\,d\#=\sum\limits_{n=1}^{\infty}f(n)$.
\item \hlt{Unit mass}: for $(X,\mathcal{M})$, $p\in X$, and $\delta_p(E)\defeq\mathbf{1}_E(p)$, $\int_Xf\,d\delta_p=f(p)$ for measurable $f:X\to[0,+\infty]$.\\
If $f:X\to\R$ is measurable, then $f\in\mathscr{L}^1(\delta_p)$ and the same formula holds.
\end{enumerate}
For $\delta_p$, verify first on indicators, then extend to simple functions by finite linearity, to nonnegative functions by monotone simple approximation, and to real-valued functions by positive--negative splitting.
\end{example}

\begin{example} \label{ex:l1-not-closed-under-squaring} \hlt{$L^1$ Is Not Closed Under Squaring}\\
On $X=\N$ with $\mathcal{M}=2^\N$, define $\mu(E)\defeq\sum\limits_{n\in E}4^{-n}$ and $f(n)\defeq\frac{2^n}{\sqrt{n}}$.\\
Then $\int_X\abs{f}\,d\mu=\sum\limits_{n=1}^{\infty}\frac{1}{2^n\sqrt{n}}<+\infty$; $\int_X f^2\,d\mu=\sum\limits_{n=1}^{\infty}\frac{1}{n}=+\infty$.\\
The first series is dominated by the geometric series $\sum2^{-n}$, while the second diverges by Corollary \ref{cor:p-series-test}.\\
Thus $f\in\mathscr{L}^1(\mu)$ need not imply $f^2\in\mathscr{L}^1(\mu)$.
\end{example}
